% GENERATED FILE. Edit canonical lessons, then run npm run generate:books.
% canonical-source-hash: fnv1a64:601783d766a4c3d3
% canonical-lessons: LA-C524-pestis, LA-C524-pestilentia, LA-C524-contagio, LA-C524-ulcus, LA-C524-tumeo

\chapter{Plague and swelling}
\label{ch:la-plague-and-swelling}
\hlchaptermodality{524}

\begin{chapteropening}
\textbf{By the end of this chapter:} \emph{I can talk about plague, infection and swelling.}

Complete the last lesson of chapter 524: I can talk about plague, infection and swelling.
\end{chapteropening}

\section[pestis, pestis]{pestis, pestis — plague, pestilence; ruin}

\label{lesson:LA-C524-pestis}



\begin{quote}
Before the new one: say the Latin for health, then the Latin for wholesome.
\end{quote}

\subsection*{You'll want to know: pestis, pestis}
\textbf{pestis, pestis} — \textquotedblleft{}plague, pestilence; ruin\textquotedblright{}.

\begin{quote}
Pestis per urbem it.
\end{quote}

\begin{quote}
\emph{Plague is going through the city.}
\end{quote}

Feminine, third declension; accusative \emph{pestem}. English \emph{pest}.

\begin{cousinweb}
\textbf{pestis} is a native Latin noun; Lewis \& Short suggests a link with \emph{perdō} \textquotedblleft{}to destroy\textquotedblright{}, but the origin is not settled. English \emph{pest} comes from it.
\end{cousinweb}

\begin{grammarlens}[title={What you've built}]
One more word for telling, explaining and describing at B1.
\end{grammarlens}

\subsection*{Your turn}
\begin{itemize}
\raggedright
  \item \emph{Say it:} \emph{pestis, pestis}
  \item \emph{Say it:} \emph{pestis, pestis}, once more
  \item \emph{Say it:} a sentence of your own with \emph{pestis, pestis}
  \item \emph{Recall:} say the Latin for health, then the Latin for wholesome, then say \emph{pestis, pestis} again
\end{itemize}

\begin{tcolorbox}[breakable,colback=teal!4,colframe=teal!35!black,title={Before you move on}]
What does pestis, pestis mean? (\textquotedblleft{}Plague, pestilence; ruin\textquotedblright{}.) Say it once more.
\end{tcolorbox}
\section[pestilentia, pestilentiae]{pestilentia, pestilentiae — an epidemic, pestilence}

\label{lesson:LA-C524-pestilentia}



\begin{quote}
Before the new one: say the Latin for wholesome, then the Latin for plague.
\end{quote}

\subsection*{You'll want to know: pestilentia, pestilentiae}
\textbf{pestilentia, pestilentiae} — \textquotedblleft{}an epidemic, pestilence\textquotedblright{}.

\begin{quote}
Eō annō pestilentia magna fuit.
\end{quote}

\begin{quote}
\emph{That year there was a great epidemic.}
\end{quote}

Feminine, first declension, from \emph{pestilēns} — plague-bearing. English \emph{pestilence}.

\begin{cousinweb}
\textbf{pestilentia} is built on \emph{pestilēns} \textquotedblleft{}plague-bearing\textquotedblright{}, from \emph{pestis} \textquotedblleft{}plague\textquotedblright{} (Lewis \& Short; EtymDB). English \emph{pestilence} comes from it.
\end{cousinweb}

\begin{grammarlens}[title={What you've built}]
One more word for telling, explaining and describing at B1.
\end{grammarlens}

\subsection*{Your turn}
\begin{itemize}
\raggedright
  \item \emph{Say it:} \emph{pestilentia, pestilentiae}
  \item \emph{Say it:} \emph{pestilentia, pestilentiae}, once more
  \item \emph{Say it:} a sentence of your own with \emph{pestilentia, pestilentiae}
  \item \emph{Recall:} say the Latin for wholesome, then the Latin for plague, then say \emph{pestilentia, pestilentiae} again
\end{itemize}

\begin{tcolorbox}[breakable,colback=teal!4,colframe=teal!35!black,title={Before you move on}]
What does pestilentia, pestilentiae mean? (\textquotedblleft{}An epidemic, pestilence\textquotedblright{}.) Say it once more.
\end{tcolorbox}
\section[contāgiō, contāgiōnis]{contāgiō, contāgiōnis — infection, contagion; touch}

\label{lesson:LA-C524-contagio}



\begin{quote}
Before the new one: say the Latin for plague, then the Latin for an epidemic.
\end{quote}

\subsection*{You'll want to know: contāgiō, contāgiōnis}
\textbf{contāgiō, contāgiōnis} — \textquotedblleft{}infection, contagion; touch\textquotedblright{}.

\begin{quote}
Morbus contāgiōne trānsit.
\end{quote}

\begin{quote}
\emph{The disease passes by contact.}
\end{quote}

Feminine, third declension, from \emph{con-} + the root of \emph{tangō} — to touch. English \emph{contagion}.

\begin{cousinweb}
\textbf{contāgiō} is built on \emph{con-} and the root of \emph{tangō} \textquotedblleft{}to touch\textquotedblright{} (Lewis \& Short): a touching together. English \emph{contagion} and \emph{contagious} come from it.
\end{cousinweb}

\begin{grammarlens}[title={What you've built}]
One more word for telling, explaining and describing at B1.
\end{grammarlens}

\subsection*{Your turn}
\begin{itemize}
\raggedright
  \item \emph{Say it:} \emph{contāgiō, contāgiōnis}
  \item \emph{Say it:} \emph{contāgiō, contāgiōnis}, once more
  \item \emph{Say it:} a sentence of your own with \emph{contāgiō, contāgiōnis}
  \item \emph{Recall:} say the Latin for plague, then the Latin for an epidemic, then say \emph{contāgiō, contāgiōnis} again
\end{itemize}

\begin{tcolorbox}[breakable,colback=teal!4,colframe=teal!35!black,title={Before you move on}]
What does contāgiō, contāgiōnis mean? (\textquotedblleft{}Infection, contagion; touch\textquotedblright{}.) Say it once more.
\end{tcolorbox}
\section[ulcus, ulceris]{ulcus, ulceris — a sore, ulcer}

\label{lesson:LA-C524-ulcus}



\begin{quote}
Before the new one: say the Latin for an epidemic, then the Latin for infection.
\end{quote}

\subsection*{You'll want to know: ulcus, ulceris}
\textbf{ulcus, ulceris} — \textquotedblleft{}a sore, ulcer\textquotedblright{}.

\begin{quote}
Medicus ulcus in crūre lavat.
\end{quote}

\begin{quote}
\emph{The doctor washes the sore on the leg.}
\end{quote}

Neuter, third declension. English \emph{ulcer}.

\begin{cousinweb}
\textbf{ulcus} is inherited from a Proto-Indo-European word for \textquotedblleft{}wound, sore\textquotedblright{}, and Greek \emph{helkos} is a cousin (EtymDB; Lewis \& Short). English \emph{ulcer} comes from it.
\end{cousinweb}

\begin{grammarlens}[title={What you've built}]
One more word for telling, explaining and describing at B1.
\end{grammarlens}

\subsection*{Your turn}
\begin{itemize}
\raggedright
  \item \emph{Say it:} \emph{ulcus, ulceris}
  \item \emph{Say it:} \emph{ulcus, ulceris}, once more
  \item \emph{Say it:} a sentence of your own with \emph{ulcus, ulceris}
  \item \emph{Recall:} say the Latin for an epidemic, then the Latin for infection, then say \emph{ulcus, ulceris} again
\end{itemize}

\begin{tcolorbox}[breakable,colback=teal!4,colframe=teal!35!black,title={Before you move on}]
What does ulcus, ulceris mean? (\textquotedblleft{}A sore, ulcer\textquotedblright{}.) Say it once more.
\end{tcolorbox}
\section[tumeō, tumēre]{tumeō, tumēre — to swell, be swollen}

\label{lesson:LA-C524-tumeo}



\begin{quote}
Before the new one: say the Latin for infection, then the Latin for a sore.
\end{quote}

\subsection*{You'll want to know: tumeō, tumēre}
\textbf{tumeō, tumēre} — \textquotedblleft{}to swell, be swollen\textquotedblright{}.

\begin{quote}
Post cāsum pēs eius tumet.
\end{quote}

\begin{quote}
\emph{After the fall his foot is swollen.}
\end{quote}

\emph{tumeō, tumēre, tumuī}. English \emph{tumour} comes from its noun \emph{tumor}.

\begin{cousinweb}
\textbf{tumeō} is inherited from a Proto-Indo-European root meaning \textquotedblleft{}to swell\textquotedblright{} (EtymDB); Lewis \& Short also compares English \emph{thumb}. English \emph{tumour} comes from its noun \emph{tumor}.
\end{cousinweb}

\begin{grammarlens}[title={What you've built}]
One more word for telling, explaining and describing at B1.
\end{grammarlens}

\subsection*{Your turn}
\begin{itemize}
\raggedright
  \item \emph{Say it:} \emph{tumeō, tumēre}
  \item \emph{Say it:} \emph{tumeō, tumēre}, once more
  \item \emph{Say it:} a sentence of your own with \emph{tumeō, tumēre}
  \item \emph{Recall:} say the Latin for infection, then the Latin for a sore, then say \emph{tumeō, tumēre} again
\end{itemize}

\begin{tcolorbox}[breakable,colback=teal!4,colframe=teal!35!black,title={Before you move on}]
What does tumeō, tumēre mean? (\textquotedblleft{}To swell, be swollen\textquotedblright{}.) Say it once more.
\end{tcolorbox}
